% !TEX program = pdflatex
\documentclass[11pt,a4paper]{article}
\usepackage[a4paper,margin=2cm,top=2.1cm,bottom=2.1cm]{geometry}
\usepackage[T5]{fontenc}
\usepackage[utf8]{inputenc}
\usepackage{mathptmx}
\usepackage{microtype}
\usepackage{booktabs,longtable,array}
\usepackage{enumitem,xcolor,hyperref,xurl,titlesec,needspace,amsmath}
\hypersetup{colorlinks=true,linkcolor=blue!45!black,urlcolor=blue!45!black,
  pdftitle={Đánh giá hiện trạng và lộ trình phát triển hệ thống MLOps cho khóa luận},
  pdfauthor={Trần Nguyễn Việt Hoàng, Bùi Ngọc Thái}}
\setlength{\parindent}{0pt}
\setlength{\parskip}{0.35em}
\setlength{\emergencystretch}{2em}
\raggedbottom
\clubpenalty=10000
\widowpenalty=10000
\setlist{leftmargin=1.5em,itemsep=0.15em,topsep=0.25em,parsep=0pt}
\titleformat{\section}{\large\bfseries}{\thesection}{0.6em}{}
\titlespacing*{\section}{0pt}{1.2ex}{0.6ex}
\titleformat{\subsection}{\normalsize\bfseries}{\thesubsection}{0.6em}{}
\titlespacing*{\subsection}{0pt}{1ex}{0.4ex}
\newcolumntype{L}[1]{>{\raggedright\arraybackslash}p{#1}}
\newcommand{\codepath}[1]{{\small\path{#1}}}
\newcommand{\statuspartial}{\textcolor{orange!65!black}{\textbf{Một phần}}}
\newcommand{\missing}{\textcolor{red!60!black}{\textbf{Chưa có}}}
\renewcommand{\tablename}{Bảng}
\begin{document}
\begin{center}
  \textbf{\Large Đánh giá hiện trạng và lộ trình phát triển\\hệ thống MLOps phục vụ khóa luận tốt nghiệp}\\[0.4em]
  Đối chiếu đề cương đã đề xuất với mã nguồn hiện tại\\
  {\small Cập nhật: 14/09/2026 -- Mốc mã nguồn: \texttt{303ceec}}
\end{center}

\section{Mục tiêu và kết luận hiện trạng}
\textbf{Đề tài:} Thiết kế và xây dựng cơ chế Continuous Training thích ứng với Data Drift trong MLOps dưới ngân sách gán nhãn hạn chế. Đề cương PDF quy định thời gian \textbf{07/09/2026--25/12/2026}, chia thành bốn giai đoạn; lộ trình này bám các mốc đó, thay cho kế hoạch 5 tháng trong bản LaTeX trước.

\textbf{Nhận định chính:} Hệ thống đã có nền tảng quản lý mô hình, suy luận, thu thập dữ liệu, báo cáo drift, training, đóng gói và triển khai. Khoảng trống lớn nhất là \textbf{vòng Continuous Training có kiểm chứng}: mở nhãn đúng ngân sách, quyết định có bằng chứng, tạo dataset tái lập, đánh giá candidate độc lập và cập nhật champion an toàn khi xảy ra lỗi bất đồng bộ.

Hai đích phải đạt song song: \textbf{RQ1} đo đánh đổi chất lượng--chi phí của kiểm chứng ít nhãn so với các trigger đối chứng; \textbf{RQ2} kiểm tra ngân sách, dữ liệu, phiên bản và thời gian phục hồi của vòng điều phối trên AWS/K3s, Argo Workflows và Kubeflow Training Operator.

\textbf{Phạm vi bằng chứng:} Đã đọc đề cương 10 trang và đối chiếu code, tests, cấu hình production cùng cây GitOps. ``Có nền'' là có implementation để kế thừa; ``Một phần'' là còn thiếu yêu cầu cụ thể; ``Chưa có'' là chưa tìm thấy implementation trong các luồng rà soát. Đây là đánh giá repo, \textbf{chưa xác nhận trạng thái cluster hoặc pipeline đã vượt kiểm thử tích hợp}; không quy đổi thành phần trăm hoàn thành.

\section{Hệ thống đã có gì và còn thiếu gì?}
\begin{longtable}{@{}L{3cm} L{1.7cm} L{\dimexpr\linewidth-4.7cm-4\tabcolsep\relax}@{}}
\caption{Đối chiếu năng lực hiện tại với đề cương.}\label{tab:gaps}\\
\toprule\textbf{Năng lực} & \textbf{Hiện trạng} & \textbf{Bằng chứng và phần cần bổ sung}\\
\midrule\endfirsthead
\toprule\textbf{Năng lực} & \textbf{Hiện trạng} & \textbf{Bằng chứng và phần cần bổ sung}\\
\midrule\endhead
Suy luận và thu thập production & \statuspartial & Có prediction ID, thời gian, features và identity theo version. Consumer ghi DB chống lặp event, kèm tín hiệu drift trong cùng transaction [E1]. Producer mới xác nhận enqueue cục bộ; cần đối soát dữ liệu thiếu trước khi chốt batch.\\
Xác suất lớp hiệu chuẩn & \statuspartial & ML worker chỉ trả max confidence 0--100, có thể null [E2]. Cần xuất/lưu $p(Y=1\mid X)$ thang 0--1, class mapping, decision threshold và calibration version; confidence hiện tại chưa đủ cho CBPE.\\
Drift và cửa sổ tái lập & \statuspartial & Có Evidently, báo cáo HTML/JSON và automatic DriftRun [E3]. Hiện lấy các dòng mới nhất theo version; cần batch snapshot S3 cố định theo ID/thời gian/hash, tách lỗi schema khỏi drift thống kê.\\
Feedback và Budget Ledger & \missing & Chưa có API mở/nhận nhãn gắn prediction, pool verify/train/gate, quota và giữ chỗ ngân sách. Đây là phụ thuộc bắt buộc của policy và replay che nhãn.\\
CBPE/ATC và kiểm chứng & \missing & Chưa có estimator, mốc sai số theo version, khoảng bất định từ nhãn thật và audit định kỳ độc lập cảnh báo.\\
Policy và MaintenanceRun & \statuspartial & Hook \texttt{handle\_drift\_detected} mới log và trả evidence [E3]; chưa quản lý lượt điều phối, bốn quyết định, yêu cầu nhãn hoặc tự submit retraining.\\
Training và Build/Register & Có nền & Có TrainingJob, snapshot, TrainingOutput và version sau build [E4--E5]. Cần tạo dataset từ production đã mở nhãn đúng pool, thay mặc định lấy data workspace mới nhất.\\
Gate và cập nhật champion & \statuspartial & Có version, metrics, alias và triển khai [E5]; chưa có gate tin cậy so sánh candidate--champion trên holdout riêng. Đổi stage/alias chưa chứng minh phiên bản phục vụ đã đổi an toàn.\\
Outbox và phục hồi CT & \statuspartial & Có EventOutbox, khóa và callback ở lifecycle hiện tại [E1, E6]. Cần nối vào MaintenanceRun và đối soát DB--Workflow--job; retry sẵn có chưa bảo đảm toàn vòng CT.\\
AWS/K3s và training production & \statuspartial & Có Terraform, Ansible, GitOps và template PyTorchJob [E7]. ConfigMap tắt training, URL training/cancel trống; overlay chưa ghi đè [E8]. Cần hoàn thiện cấu hình và smoke-test sớm.\\
Quan sát và benchmark & \statuspartial & Có metrics HTTP/Celery và ServiceMonitor nền [E9]. Chưa có dashboard ngân sách/maintenance, replay B0--B5 và bộ chèn lỗi phục vụ đề tài.\\
\bottomrule
\end{longtable}

\textbf{Chiến lược:} Hoàn thiện contract dữ liệu và ngân sách trước, đồng thời kiểm chứng đường thực thi AWS/K3s. Nối một luồng policy--training--gate--cập nhật chạy thật, rồi mở rộng thực nghiệm. Kế thừa nền tảng hiện tại để tập trung vào đóng góp của đề tài.

\section{Hướng phát triển và cách tích hợp}
\subsection{Vòng Continuous Training cần xây}
\begin{enumerate}
  \item \textbf{Suy luận $\rightarrow$ batch cố định:} ML worker xuất xác suất hiệu chuẩn; gateway phát event; Consumer lưu DB. Bổ sung kết xuất batch hợp lệ lên S3 kèm ID, version, thời gian và hash. Drift, estimator và policy dùng cùng snapshot.
  \item \textbf{Giám sát $\rightarrow$ kiểm chứng:} Kế thừa Evidently, thêm CBPE chính và ATC đối chứng. Fit/calibration và lưu mốc reference theo version. Cảnh báo chỉ gợi ý mở nhãn; audit định kỳ vẫn chạy khi cả hai tín hiệu im lặng.
  \item \textbf{Control Plane $\rightarrow$ quyết định:} Bổ sung miền CT quản lý MaintenanceRun, Feedback/LabelRequest, ledger và evidence. Policy trả KEEP, REQUEST\_LABELS, HOLD hoặc RETRAIN; tách quyết định khỏi trạng thái thực thi pending/running/failed.
  \item \textbf{RETRAIN $\rightarrow$ candidate $\rightarrow$ gate:} Tạo snapshot từ pool train production trộn dữ liệu lịch sử theo tỷ lệ cố định, lưu manifest SHA-256. Tái sử dụng TrainingJob/TrainingOutput và Build/Register; thêm pod evaluator tin cậy chạy gate. Argo điều phối execution, Kubeflow tạo job; PostgreSQL giữ trạng thái nghiệp vụ.
  \item \textbf{Gate pass $\rightarrow$ cập nhật:} Kiểm tra champion kỳ vọng, deploy candidate và xác minh health/identity trước khi chốt champion mới. Champion đã đổi thì candidate cũ bị đánh dấu superseded; gate fail/hold hoặc candidate không khởi động được phải giữ champion cũ phục vụ.
\end{enumerate}

\subsection{Contract và bất biến cần chốt}
\begin{itemize}
  \item \textbf{Snapshot và lineage:} Mỗi decision liên kết tenant/project/version/prediction/window với policy version, evidence, ledger, training job, build, deployment, gate và runtime UID. Không dùng workspace ``latest'' thay đầu vào đã chốt.
  \item \textbf{Nhãn không rò rỉ:} Tách $D_{\mathrm{train}}$, $D_{\mathrm{cal}}$, $D_{\mathrm{ref}}$ và historical holdout. Phân pool verify/train/gate bằng seed trước khi đọc nhãn, không trùng ID hoặc dùng nhãn chéo trong thực nghiệm chính. Nhãn ẩn/tương lai chỉ thuộc evaluator ngoại tuyến.
  \item \textbf{Ngân sách nguyên tử:} Khóa và ràng buộc DB bảo đảm
  \[
    B_{\mathrm{used}}=B_{\mathrm{verify}}+B_{\mathrm{train}}+B_{\mathrm{gate}},
    \qquad B_{\mathrm{used}}+B_{\mathrm{reserved}}\leq B.
  \]
  Retry cùng ID không chi lại; mọi nhãn đã mở vẫn tính phí dù không retrain hoặc candidate bị loại. Chỉ giải phóng quota giữ chưa tiêu sau khi đối soát kết quả.
  \item \textbf{Kiểm chứng có bất định:} Khoảng $[L_t,U_t]$ của mức tăng lỗi phải xét sai số cả production và reference. $L_t>\delta_{\mathrm{harm}}$ xác nhận suy giảm; $U_t\leq\delta_{\mathrm{keep}}$ hỗ trợ KEEP. Tối đa hai lượt lấy nhãn định trước, kiểm soát việc xem kết quả lặp. Thiếu bằng chứng/quota/dữ liệu hoặc batch lỗi thì HOLD có lý do và lịch xét lại.
  \item \textbf{Gate độc lập:} So sánh ghép cặp candidate--champion trên pool gate riêng đã tính phí; chỉ pass khi lợi ích vượt biên định trước và không suy giảm quá biên trên historical holdout. Calibration của candidate độc lập training. Metrics do training gửi lên không thay evaluator tin cậy.
  \item \textbf{Phục hồi:} Lưu trạng thái và outbox cùng transaction; dispatch sau commit. Dùng khóa nghiệp vụ, attempt ID, callback xác thực và worker đối soát UID trước retry. Callback trễ không được hồi sinh trạng thái kết thúc hoặc cập nhật từ champion cũ.
\end{itemize}

\textbf{Giữ phạm vi gọn:} Hai dataset tabular nhị phân, một model family có xác suất hiệu chuẩn, một recipe full-retraining cố định và CPU cho MVP. Không mở rộng sang pseudo-labeling, incremental learning, tuning động, chẩn đoán nhân quả, đa GPU hay HA đa vùng. UI nâng cao chỉ thực hiện sau các đầu ra bắt buộc.

\section{Lộ trình đến ngày 25/12/2026}
Các mốc theo Mục VI của PDF; tiêu chí dưới đây dùng để kiểm tra tiến độ. Ngày 14/09 đang ở tuần 2 theo lịch, \textbf{không đồng nghĩa các việc GĐ1 đã hoàn thành}.

\begin{longtable}{@{}L{2.3cm} L{\dimexpr(\linewidth-2.3cm-4\tabcolsep)/2\relax} L{\dimexpr(\linewidth-2.3cm-4\tabcolsep)/2\relax}@{}}
\caption{Bốn giai đoạn và đầu ra kiểm chứng được.}\\
\toprule\textbf{Thời gian} & \textbf{Công việc trọng tâm} & \textbf{Đầu ra/tiêu chí chuyển giai đoạn}\\
\midrule\endfirsthead
\toprule\textbf{Thời gian} & \textbf{Công việc trọng tâm} & \textbf{Đầu ra/tiêu chí chuyển giai đoạn}\\
\midrule\endhead
GĐ1 -- Tuần 1--4\newline 07/09--04/10 & Chọn dataset/model/recipe; pilot batch, calibration, quota. Chốt schema, trạng thái, gate và lineage. Rà cấu hình cloud, training/webhook; smoke-test Argo/Kubeflow. & Protocol v1 và contract thống nhất. Một job CPU từ submit đến output/callback thật, có UID/log và metrics; chốt cấu hình môi trường, quota và trần chi phí AWS.\\
GĐ2 -- Tuần 5--8\newline 05/10--01/11 & Chuẩn hóa probability/log; window/S3 manifest; API nhãn, pool và ledger. Tích hợp CBPE/ATC, kiểm chứng và replay che nhãn. Dựng khung policy/baseline để kiểm thử. & Replay cùng seed cho cùng kết quả; không đọc nhãn ẩn/trùng pool. Tests local kiểm tra quota đồng thời, nhãn trùng, snapshot và estimator theo version.\\
GĐ3 -- Tuần 9--12\newline 02/11--29/11 & Hoàn thiện policy và B0--B5. Nối snapshot--training--build--gate--cập nhật trên AWS/K3s; outbox CT, reconciliation, dashboard/cảnh báo. & Một vòng cập nhật thật có lineage đầy đủ; gate fail/hold giữ champion. Retry/callback trùng không chi lại hoặc tạo logical job trùng; dashboard nhận metrics thật.\\
GĐ4 -- Tuần 13--16\newline 30/11--25/12 & Chạy RQ1 hai dataset và ablation; chèn lỗi RQ2; phân tích chất lượng/chi phí/phục hồi. Hoàn thiện báo cáo, demo và slides. & Kết quả B0--B5, chi phí nhãn/compute và toàn bộ nhóm lỗi có trace. Có script, seed/manifest, cấu hình tái lập và báo cáo giới hạn.\\
\bottomrule
\end{longtable}

\textbf{Cách tổ chức:} Đề cương đặt hoàn thiện policy ở GĐ3; cần chốt đặc tả từ GĐ1 và dựng khung ở GĐ2 để replay, ledger và execution dùng cùng contract. Viết báo cáo/lưu bằng chứng theo từng mốc; dành 14/12--25/12 chủ yếu cho sửa lỗi, phân tích và bàn giao.

\textbf{Phân công:} Hoàng phụ trách dữ liệu, estimator, policy, replay và RQ1; Thái phụ trách API/ledger, cloud, Argo/outbox/reconciliation, metrics và RQ2. Cả hai cùng chốt contract, gate, cập nhật và nghiệm thu end-to-end; bàn giao bằng schema, fixture và tiêu chí kiểm tra chung.

\subsection{Việc cần làm ngay trong phần còn lại của GĐ1}
\begin{samepage}
\begin{enumerate}
  \item Chốt hai dataset, model family và split; pilot xem quota có đủ mẫu/lớp cho verify, train và gate.
  \item Viết contract v1 cho probability event, window manifest, API nhãn/ledger, MaintenanceRun và gate; thống nhất ownership và liên kết ID.
  \item Xác minh lý do training đang tắt; chuẩn bị cấu hình môi trường thực nghiệm và URL training/cancel, rồi chạy job CPU nhỏ kiểm tra Argo--Kubeflow--S3--callback sau khi prerequisites đạt.
  \item Tạo issue theo phụ thuộc: probability/window $\rightarrow$ feedback/ledger $\rightarrow$ verification/policy $\rightarrow$ snapshot/gate $\rightarrow$ cập nhật/phục hồi $\rightarrow$ benchmark. Cập nhật roadmap bằng bằng chứng sau mỗi mốc.
\end{enumerate}
\end{samepage}

\section{Đánh giá và điều kiện hoàn thành}
\subsection{RQ1: chất lượng dưới ngân sách hữu hạn}
So sánh \textbf{B0} giữ nguyên model; \textbf{B1} trigger theo drift; \textbf{B2} trigger theo CBPE/ATC; \textbf{B3} audit định kỳ; \textbf{B4} policy đề xuất; \textbf{B5} giám sát đầy đủ nhãn. B1--B4 dùng cùng trần ngân sách, lịch/pool nhãn training, recipe và gate; B0 không cập nhật, B5 báo riêng lợi thế thông tin và chi phí, không coi là baseline cùng ngân sách.

Bao phủ sáu kịch bản PDF: ổn định, covariate shift, phân phối quay lại, concept drift/mất calibration, feedback thiên lệch và không có nhãn bổ sung. Làm ba kịch bản đầu trước; ba kịch bản sau kiểm tra giới hạn nhưng vẫn cần kết quả theo phạm vi đề cương. Không có nhãn thì đo giám sát/HOLD, không tuyên bố giải quyết supervised retraining.

\textbf{Chỉ số:} Error/F1/recall, MAE estimator, false trigger/missed degradation, thời gian phục hồi chất lượng; nhãn verify/train/gate, số job/candidate bị loại, thời gian và CPU-hours. Khóa ngưỡng, seed, số lần lặp, batch và quota trên pilot/development trước test; so sánh theo cặp ở cấp episode/block thời gian, báo khoảng bất định. Không hứa mức tiết kiệm hoặc tăng accuracy trước thực nghiệm.

\subsection{RQ2: bất biến và phục hồi trên AWS/K3s}
Kiểm tra sáu nhóm lỗi: nhãn trùng/request đồng thời gần hết quota; worker/pod dừng sau commit/dispatch; callback trễ/lặp/đảo thứ tự; training/build/gate lỗi/timeout; workspace/S3 hoặc champion đổi giữa các bước; gate không đạt hoặc candidate crash-loop. Lưu điểm chèn lỗi, attempt, runtime UID và kết quả đối soát; so sánh với lượt không lỗi cùng cấu hình.

\textbf{Nghiệm thu theo PDF:} Vòng CT chạy thật từ giám sát đến cập nhật; vượt toàn bộ bài kiểm thử bất biến đã xác định; \textbf{0 vi phạm ngân sách, 0 logical job/cập nhật trùng, 0 cập nhật sai phiên bản và 100\% decision có lineage}. Đo median/p95 thời gian từ lúc thành phần lỗi hoạt động lại đến khi luồng tiếp tục hoặc kết thúc có kiểm soát, kèm số lần lặp. Kết quả chỉ áp dụng cho tập kiểm thử hữu hạn.

\textbf{Bàn giao:} Mã nguồn/tests; Terraform/Ansible/GitOps tái lập; workflow training/gate/cập nhật; dashboard/cảnh báo CT có metrics thật; replay, fault harness, seed/manifest và kết quả; báo cáo RQ1/RQ2 và demo. Docker local không thay bằng chứng AWS/K3s.

\clearpage
\section{Rủi ro và căn cứ rà soát}
\begin{itemize}
  \item \textbf{Nhãn ít/calibration sai:} Pilot sớm, audit độc lập cảnh báo, HOLD khi khoảng còn rộng; báo riêng giới hạn estimator và mẫu nhãn thiên lệch.
  \item \textbf{Cloud chậm hoặc tốn chi phí:} Smoke-test GĐ1, CPU/dataset nhỏ cho RQ2, gom lịch chạy và tách chi phí nền/job. Nếu cần thu hẹp cấu hình/lần lặp, thống nhất với GVHD; không tự bỏ yêu cầu bắt buộc.
  \item \textbf{Tích hợp muộn/vượt phạm vi:} Làm một vòng CT hoàn chỉnh trước benchmark, kiểm tra ledger/retry/gate ngay khi tích hợp; cắt UI nâng cao và model family bổ sung trước đầu ra đề cương.
\end{itemize}

\textbf{Nguồn yêu cầu:} \codepath{references/DeCuongKLTN_HoangTNV_ThaiBN.pdf} (Mục II, IV, V, VI). Mã E trong Bảng~\ref{tab:gaps} dẫn tới các file sau, tính từ gốc repo. Các khối mới ở Mục 3 là đề xuất thiết kế, chưa phải implementation đã có.

\textbf{Quy ước đường dẫn:} Trong E3--E6, các đường dẫn bắt đầu bằng \codepath{apps/} có tiền tố \codepath{services/control-plane/src/} được lược bỏ.
\begingroup\footnotesize
\setlength{\parskip}{0pt}
\renewcommand{\codepath}[1]{\path{#1}}
\begin{description}[font=\normalfont\bfseries,leftmargin=2.5em,labelwidth=2em,itemsep=0pt,topsep=0pt]
  \item[E1] \codepath{services/model-server/src/api.py}; \codepath{services/consumer/src/kafka_runtime.py}, \codepath{services/consumer/src/database.py}: event, persistence, offset và automatic-drift outbox.
  \item[E2] \codepath{services/machine-learning-serving/src/api.py}: predict và max confidence.
  \item[E3] \codepath{services/evidently/src/main.py}; \codepath{apps/drift/services/automatic.py}; \codepath{apps/drift/tasks.py}: drift, watermark và hook CT.
  \item[E4] \codepath{apps/training/services/jobs.py}; \codepath{apps/training/tests/test_argo_backend.py}: snapshot/submit, tests backend dùng dependency mô phỏng.
  \item[E5] \codepath{apps/registry/models.py}; \codepath{apps/registry/services/versions.py}; \codepath{apps/registry/services/routing.py}: version, metrics, alias và endpoint healthy.
  \item[E6] \codepath{apps/observability/services/outbox.py}; \codepath{apps/training/api/webhooks.py}: outbox và callback hiện tại.
  \item[E7] \codepath{infra/main.tf}; \codepath{ansible/site.yml}; \codepath{k8s/gitops/production/}; \codepath{k8s/argo/workflows/training-workflowtemplate.yaml}: IaC/GitOps và training qua PyTorchJob.
  \item[E8] \codepath{k8s/apps/base/control-plane/configmap.yaml}; \codepath{k8s/apps/overlays/production/control-plane/kustomization.yaml}: training tắt, URL training/cancel trống, overlay không bổ sung giá trị.
  \item[E9] \codepath{services/control-plane/src/common/metrics.py}; \codepath{k8s/platform/observability/kustomization.yaml}: metrics và quan sát nền.
\end{description}
\endgroup

\end{document}
